\section{Categorització del sexisme en continguts multimèdia}\label{ch:marc}

L'estudi de \textcite{mused} estableix els fonaments conceptuals per a la classificació del sexisme en continguts multimèdia de les xarxes socials.
Defineixen el sexisme com ``prejudici i discriminació basats en el sexe o el gènere'', en un sentit ampli, que inclou l'orientació sexual i la identitat de gènere.
Aquesta definició segueix l'enfocament d'\textcite{o2009encyclopedia}, que conceptualitza el sexisme com una ideologia que reprodueix i marginalitza subjectivitats que no s'ajusten a les normes heteropatriarcals.

Dins d'aquest marc conceptual, \textcite{mused} identifiquen diversos reptes per a la classificació del sexisme:
\begin{itemize}
    \item La distinció entre sexisme implícit i explícit. Les formes implícites, com els estereotips subtils, són més subjectives, depenen del context cultural i obtenen un acord menor entre els anotadors.
    \item La importància del context multimodal. Alguns tipus de sexisme, com l'objectificació, només es poden detectar adequadament si s'analitzen diverses modalitats, i no únicament el text. Això pot generar contradiccions quan un missatge textual neutre s'acompanya d'elements visuals sexistes.
    \item La dependència cultural del fenomen, que pot provocar una variabilitat notable en els patrons d'atenció dels participants.
\end{itemize}
Aquests factors tenen implicacions directes en els patrons d'atenció entre humans i models d'intel·ligència artificial, especialment en categories de sexisme subtils o contextuals, i motiven l'ús de tècniques d'eye-tracking i d'interpretabilitat computacional que desenvolupem a continuació.

\section{Corpus d'origen: MuSeD}\label{sec:mused}

El corpus utilitzat en aquest treball prové de MuSeD (\textit{Multimodal Sexism Detection Dataset})~\cite{mused}, un conjunt de dades per a la detecció de sexisme en vídeos en castellà.
MuSeD consta de 400 vídeos extrets de dues xarxes socials: TikTok i BitChute.
TikTok, tot i prohibir explícitament la misogínia a les seves directrius comunitàries, continua allotjant contingut sexista de manera explícita i implícita; BitChute, en canvi, és una plataforma de baixa moderació.
El conjunt de dades és aproximadament equilibrat, amb un 48,5\% de vídeos etiquetats com a sexistes i un 51,5\% com a no sexistes.

\subsection{Recollida de dades}

Els vídeos es van recollir entre maig i octubre de 2024.
%Per a TikTok, 
Es va elaborar una llista de 187 \textit{hashtags} en castellà que cobria temàtiques relacionades amb el gènere, la sexualitat, el feminisme i la discriminació per sexe o orientació sexual. 
%(per exemple, \textit{\#brechasalarial}, \textit{\#violenciamachista} o \textit{\#lgbt}).
%Per a BitChute, es van adaptar alguns d'aquests \textit{hashtags} com a paraules clau de cerca.
Aproximadament el 90\% dels vídeos provenen de TikTok. 
Les transcripcions textuals es van obtenir automàticament amb Whisper~\cite{whisper} i van ser revisades i corregides per una lingüista professional.

\subsection{Procés d'anotació}

L'anotació es va dur a terme en tres nivells: text (transcripció i OCR\footnote{Optical character recognition.}---text escrit al vídeo), àudio i vídeo complet.
Sis anotadors amb formació en lingüística, dividits en dos equips, van classificar cada ítem com a sexista o no sexista.
Addicionalment, van etiquetar els segments (\textit{spans}) textuals o segments de vídeo que motivaven la classificació.
El marc d'anotació distingia quatre tipus de contingut sexista:
\begin{itemize}
    \item \textbf{Estereotip} (\textit{Stereotype}): afirmacions sobre propietats que suposadament diferencien homes i dones a partir de creences estereotipades~\cite{samory2021call}. Reforça rols de gènere tradicionals i presenta diferències com a ``naturals'' o biològicament determinades.
    \item \textbf{Desigualtat} (\textit{Inequality}): discurs que nega
    l'existència de desigualtats de gènere o presenta el moviment feminista com una amenaça per als homes. Inclou la victimització masculina en polítiques d'igualtat.
    \item \textbf{Discriminació} (\textit{Discrimination}): contingut que
    discrimina per orientació sexual o identitat de gènere i que denigra la comunitat LGBTQ+~\cite{chakravarthi2021dataset}.
    \item \textbf{Objectificació} (\textit{Objectification}): representació de les dones com a objectes físics valorats principalment per la seva utilitat~\cite{szymanski2011sexual}.
\end{itemize}

Entre els vídeos sexistes, el 56,2\% correspon a estereotips, el 39,2\% a desigualtat, el 15,5\% a discriminació i el 3,6\% a objectificació.
Com que un mateix ítem pot rebre més d'una categoria, aquests percentatges no han de sumar el 100\%.
Cal destacar que, si un ítem conté una crítica, una denúncia o el relat d'una experiència sexista, s'etiqueta com a no sexista.
A més, també es va anotar la implicitud del sexisme~\cite{schmeisser2022criteria}, així com la presència d'humor o ironia.


\subsection{Estructura del corpus}

Cada instància del corpus proporciona les variables següents, rellevants per als criteris de selecció descrits a la \autoref{sec:filtratge}:
\begin{itemize}
    \item \texttt{sexist\_text}: etiqueta binària que indica si el text, aïllat del context audiovisual, es considera sexista.
    \item \texttt{sexist\_video}: etiqueta binària equivalent, però basada en el contingut multimodal complet (vídeo, àudio i text).
    \item \texttt{sexist\_soft}: valor entre 0 i 1 que reflecteix el grau d'acord entre els anotadors. Té quatre valors possibles: 0, 0,33, 0,67 i 1.
    \item Anotacions a escala de segments: fragments concrets del text marcats com a sexistes per part dels anotadors, que permeten identificar quines parts del text motiven l'etiqueta global.
\end{itemize}

\subsection{Acord entre anotadors}

La fiabilitat de les anotacions es va mesurar amb la $\kappa$ de Fleiss~\cite{fleiss1971measuring}.
En l'anotació textual, es van obtenir valors de $\kappa=0,719$ (equip~1) i $0,717$ (equip~2). En l'anotació del vídeo complet, els valors van augmentar fins a $0,834$ i $0,853$, respectivament, cosa que indica un acord \textit{substancial}~\cite{landis1977measurement}.
La millora associada a la incorporació de modalitats addicionals suggereix que, en alguns casos, el sexisme és difícil de detectar només a partir del text. Això pot deure's a la dificultat d'identificar formes implícites o a la contradicció entre el contingut verbal i el no verbal.
Precisament aquest fet motiva el criteri de coherència entre modalitats descrit a la metodologia (\autoref{sec:filtratge}).


\section{Lectura i eye-tracking}

La lectura implica una sèrie de moviments oculars: fixacions---períodes de 200--300\,ms en què l'ull està relativament estàtic---intercalades amb sacades ràpides, que duren aproximadament 30--50\,ms~\cite{rayner1998eye}.
Durant les fixacions, el lector processa la informació d'una àrea focal de 3--4 lletres; la informació perifèrica guia les sacades successives cap a les paraules següents.
L'eye-tracking és una tècnica no invasiva que permet registrar els moviments oculars durant la lectura.
Ens aporta informació indirecta, però objectiva, sobre els processos cognitius \cite{conklinEyeTrackingGuideApplied2018}.
Aquesta idea s'anomena \textit{mind-eye hypothesis} (hipòtesi ull-ment)~\cite{just1980theory}.
La hipòtesi es basa en dues premisses~\cite{pickering2004eye}: (1) allò en què es fixen els ulls s'està processant mentalment i (2) com més temps es dedica a una regió, més esforç cognitiu requereix.

Així doncs, els moviments oculars proporcionen una mesura en temps real del processament lingüístic: les fixacions reflecteixen l'atenció i l'esforç cognitiu, mentre que les sacades indiquen els moviments entre diferents regions del text.
En la literatura sobre eye-tracking, s'anomena àrea o regió d'interès cada zona de la imatge que es vol estudiar.
En els experiments de lectura, les àrees d'interès poden correspondre a caràcters, paraules, sintagmes o fragments de text.
En aquest treball, les àrees d'interès corresponen a paraules.
Les mètriques d'eye-tracking més utilitzades en la recerca sobre lectura inclouen:

\begin{itemize}
    \item \textbf{Durada total de fixació (TFD, \textit{Total Fixation Duration})}: suma de les durades de les fixacions sobre una paraula al llarg de tota la lectura, incloses les relectures. Capta l'atenció acumulada i el processament deliberat, i és especialment sensible a la dificultat semàntica i pragmàtica del text. En aquest treball és la mètrica principal per comparar l'atenció humana amb les distribucions de prominència dels models, ja que totes dues representen l'atenció acumulada sobre els tokens. La TFD es normalitza pel temps total de fixació de cada text.
    \item \textbf{Durada de la primera fixació (FFD, \textit{First Fixation Duration})}: durada de la primera fixació sobre una paraula. Capta la reacció inicial al contingut i reflecteix el processament superficial i el reconeixement lèxic ràpid~\cite{rayner2009eye}. Una FFD més llarga pot indicar que una paraula és menys predictible o més complexa cognitivament.
    \item \textbf{Nombre de fixacions (FC, \textit{Fixation Count})}: nombre total de fixacions sobre una paraula. Reflecteix l'atenció acumulada, incloses les relectures i el processament deliberat~\cite{rayner2009eye}.
    \item \textbf{Regressions}: moviments oculars cap enrere (dreta a esquerra en escrits d'esquerra a dreta) que reflecteixen reprocessament, verificació o resolució d'ambigüitat. Taxes més altes de regressió indiquen fragments del text que generen més dificultat de processament.
\end{itemize}

Aquestes mètriques informen no només del processament superficial, com el reconeixement de paraules, sinó també de la integració semàntica i pragmàtica. Per a més informació sobre la TFD i les regressions, vegeu l'\autoref{app:tfd}.

\subsection{Corpus d'eye-tracking per al PLN}

En el camp del PLN, existeixen diversos corpus d'eye-tracking per entrenar i avaluar models computacionals.
Aquests corpus permeten entrenar models predictius de fixacions i explorar la relació entre l'atenció humana i els mecanismes d'atenció dels models de llenguatge.

ZuCo \cite{hollenstein_zuco_2018} és un recurs simultani d'electroencefalografia (EEG) i eye-tracking amb 12 subjectes que llegeixen frases angleses.
El corpus inclou 21.629 paraules en 1.107 frases, amb 154.173 fixacions enregistrades.
Els materials inclouen frases de ressenyes de pel·lícules amb etiquetes de sentiment i paràgrafs biogràfics de Wikipedia amb relacions semàntiques anotades.
Els participants van executar tres tasques: dues de lectura normal, amb materials diferents, i una tasca específica de cerca de relacions.
Els autors van extreure diferents característiques oculars per a cada paraula i frase, com la FFD o el temps total de lectura.

ZuCo 2.0 \cite{hollenstein_zuco_2020} amplia el predecessor amb 18 subjectes i contrasta la lectura normal amb la lectura específica per la tasca. Demostra que l'objectiu de la lectura modifica profundament els patrons d'atenció visual, una qüestió que també abordem quan comparem l'atenció humana (lectura per comprensió) amb les atribucions del model (classificació de sexisme).

\textcite{arcosMultimodalSensorData2024} aplica dades d'eye-tracking, freqüència cardíaca i EEG a la detecció de sexisme en mems.
Mostren que els mems sexistes requereixen més esforç cognitiu, que es manifesta en temps de reacció més llargs i un nombre de fixacions superior. Aquest efecte augmenta progressivament dels mems no sexistes als mems sexistes crítics.
A més, el seu model, que integra text i imatge amb EEG i eye-tracking, aconsegueix una AUC\footnote{L'AUC (\textit{Area Under the Curve}) representa l'àrea sota la corba ROC, que relaciona la taxa de veritables positius amb la taxa de falsos positius.} de 0,8 en la detecció binària. Això suposa una millora del 3,4\% respecte d'un model base que només utilitza text i imatge.
Per tant, les respostes fisiològiques proporcionen un senyal complementari al contingut, i la fusió multimodal millora el rendiment en la detecció de sexisme.


\section{Comparació entre atenció humana i computacional}

La intersecció entre els patrons d'atenció humans i els mecanismes d'atenció computacionals s'ha convertit en un tema rellevant per a la interpretabilitat~\cite{lipton2018mythos} de la IA i per a la lingüística cognitiva.
\textcite{soodInterpretingAttentionModels2020} van proposar un mètode sistemàtic per comparar l'atenció humana i la computacional en tasques de comprensió lectora. Per fer-ho, van utilitzar el corpus d'eye-tracking MQA-RC\footnote{MovieQA Reading Comprehension, QA=Question Answering}, amb 23 subjectes.
Van observar que, en els models basats en LSTM i CNN, hi havia una correlació significativa entre la similitud amb l'atenció humana i el rendiment (Spearman $\rho = -0,73$ i $-0,72$, respectivament, $p < 0,001$). Aquesta relació, però, no es mantenia en models Transformer com XLNet~\cite{yang2019xlnet} ($\rho = -0,16$, $p = 0,381$).
A més, la divergència de Kullback-Leibler (KL) entre l'atenció LSTM i humana ($0,018$) és significativament menor que la de XLNet ($0,022$), fet que suggereix que diferents arquitectures aprenen estratègies d'atenció diferents i que la similitud amb l'atenció humana no garanteix el millor rendiment.
Aquesta observació té implicacions per al nostre treball: si els models Transformer no mostren una correlació significativa amb l'atenció humana, cal explorar també altres fonts d'informació, com les anotacions humanes i les dades d'eye-tracking.

\textcite{ikhwantriLookingDeepEyes2023} estableixen el marc metodològic per a la comparació sistemàtica entre patrons d'atenció humans i computacionals en tasques de PLN.
Analitzen tres tasques (anàlisi de sentiment, classificació de relacions, pregunta-resposta), quatre mètodes d'interpretabilitat (gradient simple, Integrated Gradients, pertorbació d'entrada i atenció) i tres arquitectures (LSTM, CNN i Transformer).
Es basen en dos conjunts de dades d'eye-tracking: ZuCo i MQA-RC.

Per mesurar la conformitat entre la prominència del model i l'atenció humana, proposen calcular la divergència KL entre la distribució de prominència sobre les paraules d'entrada---obtinguda amb un mètode d'interpretabilitat---i una característica d'eye-tracking (FC o FFD).
Una KL més petita indica una similitud més gran entre les distribucions: el model assigna més prominència a les mateixes paraules que les persones.
Els seus resultats principals són:

\begin{itemize}
    \item La KL és consistentment més petita amb FFD que amb FC en gairebé totes les combinacions de tasques i mètodes.
    % Això s'explica perquè la FFD captura la primera lectura d'una paraula, similar al comportament dels models que processen el text una sola vegada, mentre que el FC inclou relectures que els models no realitzen. % No és un bon anàleg. Un model no rellegeix.
    \item Els mètodes basats en gradients obtenen una KL menor que els mètodes de pertorbació (\textit{input perturbation}, IP) i d'atenció (Attn) en les tasques d'anàlisi de sentiment i classificació de relacions. En canvi, en la tasca de QA amb resposta múltiple, basada en texts més llargs, l'atenció és el mètode amb la KL més baixa.
    \item El Transformer mostra la KL més baixa en la majoria de les combinacions, excepte amb el mètode IP en l'anàlisi de sentiment i amb Attn en la classificació de relacions.
    \item En anàlisi de sentiment, els humans centren l'atenció en noms propis, adjectius, noms, adverbis i nombres; els models amb el millor mètode (Grad-Transformer) sobrevaloren els adjectius, però subvaloren els noms propis i els nombres. En classificació de relacions, els models no capten bé els verbs i adjectius que indiquen relacions entre entitats.
    \item Les corbes KL--Performance que proposen tendeixen a ser decreixents en tasques de classificació (més conformitat implica més rendiment), mentre que en QA són més planes o irregulars.
\end{itemize}

Els resultats d'\textcite{ikhwantriLookingDeepEyes2023} suggereixen que la relació entre conformitat humana i rendiment varia segons la tasca, el mètode d'interpretabilitat i l'arquitectura.
En el nostre cas, la classificació de texts és similar a les tasques de sentiment i de relacions.
Per això, esperem que els mètodes basats en gradients mostrin un alineament més gran amb l'atenció humana.
A més, la nostra aproximació tripartida (anotacions, eye-tracking i models) ens permet abordar la necessitat d'integrar diverses fonts d'informació per entendre el comportament dels models.

\subsection{Variabilitat en explicacions: humans vs.\ models}

\textcite{bogaert_explanation_2026} comparen anotacions humanes, basades en tokens ressaltats, amb explicacions generades mitjançant Layer-wise Relevance Propagation (vegeu la \autoref{sec:metodes}) per models transformer adaptats.
Identifiquen diferències notables en diversos àmbits.
Pel que fa a la granularitat, els humans destaquen segments més llargs i semànticament rellevants, mentre que els models produeixen atribucions més difuses a escala de token.
Les persones mostren més acord en els tokens més rellevants, mentre que els models tendeixen a coincidir més quan es consideren tots els tokens.
Pel que fa a la sensibilitat al tipus de contingut, els humans mostren més acord en texts d'opinió, mentre que els models són més consistents en texts de notícies.

Aquesta variabilitat en les explicacions no és exclusiva de la comparació entre persones i models; també es manifesta entre models funcionalment equivalents.
L'efecte Rashomon~\cite{breiman2003statistical} estableix que els models amb un rendiment similar poden produir explicacions radicalment diferents per a una mateixa predicció. Això posa en dubte la noció de \textit{l'única} explicació correcta d'un model.
Si els models poden arribar a la mateixa conclusió per camins diferents, l'alineament amb els processos humans es converteix en un criteri de disseny explícit. En aquest context, la variabilitat de les explicacions és una propietat que cal quantificar, no pas un soroll que cal eliminar.

\subsection{Plausibilitat i fidelitat}

Adoptem el marc de \textcite{jacovi2020towards}, que distingeix entre plausibilitat i fidelitat.
La plausibilitat és la mesura en què una explicació resulta convincent per a un observador humà i s'alinea amb les seves expectatives sobre els factors rellevants.
La fidelitat és la mesura en què reflecteix el procés real de decisió del model.

L'enfocament tripartit que proposem permet situar cadascuna de les fonts de dades en aquest marc.
Les anotacions humanes, basades en segments ressaltats, permeten avaluar la plausibilitat de les explicacions.
Les dades d'eye-tracking ofereixen una mesura més objectiva dels processos d'atenció, tot i que no constitueixen una mesura directa de la fidelitat del model.
La combinació de totes dues fonts permet explorar la possible tensió entre plausibilitat i fidelitat.

\section{Enfocament tripartit i hipòtesis del treball}

El nostre treball amplia el marc de~\textcite{bogaert_explanation_2026} incorporant dades d'eye-tracking, que capturen els patrons d'atenció visual durant la lectura.
Aquesta aproximació permet establir tres comparacions: entre les atribucions dels models i els segments textuals que els anotadors humans han marcat com a rellevants; entre les atribucions dels models i les mètriques d'eye-tracking per token o regió textual; i entre els patrons d'atenció visual capturats per eye-tracking i les decisions conscients de ressaltar fragments rellevants.

A partir d'aquest marc, plantegem tres hipòtesis:

\begin{itemize}
    \item \textbf{H1}: Els models adaptats mostraran un alineament moderat amb els patrons humans, amb diferències notables de granularitat, seguint~\textcite{bogaert_explanation_2026}.
    \item \textbf{H2}: L'alineament serà major per a sexisme explícit que per a casos de sexisme implícit, humor o ironia.
    \item \textbf{H3}: Les tres fonts de dades proporcionaran informació complementària entre si.
\end{itemize}
